\section{The two missing weak-IH source endpoints}
\label{sec:weak-ih-endpoints}
We use the exact parent, graph, roots, coefficients $K,k$, and both
normal factorizations (IC3)--(IC12) in the preserved interior-corner
proof.  Those identities are exact for every positive canonical C2
source; their validity does not require the interior-shape assumption.
Here $z$ and testing $u,w$ are infinitesimal (zero testing coordinates
are allowed), $B$ is unbounded, $H\in(0,1)$, $\rho=B/y$, and
$N=4n+2$, $\nu=Nr$. All source and testing rates are independent.

\paragraph{The endpoint $\rho_0=0$.}
Put $s=r/\rho$ and $\tau=1/B$. Nonnegativity of $n$ and the exact
source formulas imply that $s$ is bounded. The following normalized
source quantities have the indicated residues:
\[
 \nu/\rho\to2,\quad nr/\rho\to(1-s)/2,\quad
 D/\rho^2\to1/2,\quad BC,Bm/z^2,AH/B,\alpha(1-H)/B\to1/2,
\]
\[
 k/(B\rho)\to2s,\qquad 0\le s_0\le1.
\tag{WIE1}
\]
These are uniform opening identities: after substituting
$B=1/\tau$, $y=1/(\tau\rho)$, $r=s\rho$, their normalized rational
functions have a nonzero denominator at $z=\tau=\rho=0$, including
$H=0,1$. The fresh checker records these denominator residues.
Thus they concern simultaneous, not iterated, infinitesimal limits.

Replace $T$ by $T_\rho=\rho T$ in (IC7). Its coefficients are bounded;
for example $r^2n/\rho^2=s(nr/\rho)$ is bounded. Add the exact zero
\[
 T_\rho^2+(u^2-z)T_\rho R_2
 -\rho u(zR_0-uR_1)(uR_0-R_1)=0.
\tag{WIE2}
\]
This is (IC8a) multiplied by $\rho$, and uses no mixed inverse.
The residue Gram has $R_0,R_1,E$ diagonal entries $2,2,2$, and
\[
 (R_2,T_\rho):\quad
 \begin{pmatrix}2&-s_0\\-s_0&1+2s_0\end{pmatrix}.
\]
Its determinant $2+4s_0-s_0^2\ge2$ is strictly positive on $[0,1]$.

The entire normal restoration survives this rescaling. In (IC11),
replace $f_XT$ by $(f_X/\rho)T_\rho$, where
\[
 f_X/\rho=-2(nr/\rho)(1+w^2)-s(2D+1)
\]
is bounded. Every multiplier still contains testing $u$ or $w$, and
$1/\sqrt N$ is bounded. The remaining coefficients in (IC11) are
unchanged and infinitesimal. Formula (IC12) uses the root $R_1$ rather
than $T$ and remains bounded and infinitesimal unchanged. The normal
quadratic residue is $\operatorname{diag}(1,2)$. Hence the full
seven-root Gram is strict. This includes arbitrary positive $r,H,1-H$
and literal $n=0$.

\paragraph{The endpoint $\rho_0=1$.}
Set $e=1+zr-\rho+H/y=j>0$, an infinitesimal, and write
\[
 \rho=\frac{1+zr-e}{1-H\tau},\qquad \tau=1/B.
\]
The exact source formulas give the simultaneous normalized residues
\[
 De\to1/4,\quad\nu\to1,\quad nr\to(1-2r)/4,\quad
 BC\to1/4,\quad Bm/z^2\to(1-r)/4,
\]
\[
 AH/B,\alpha(1-H)/B\to1/4,\quad
 0\le r_0\le1/2.
\tag{WIE3}
\]
Put $\kappa=k-8DHr$; then $\kappa/B\to r$.
The scaled rational functions have regular nonzero denominators at
$z=\tau=e=0$ after the displayed powers of $e$ are absorbed.

The large coefficient $D$ must be retained in a square, not treated
as a bounded Gram coefficient. Introduce the full-state root
\[
 F_D=\sqrt D\,(J+eZ),\qquad a=\sqrt{H/B},\qquad
 Q=R_2-rT-aE.
\]
On the graph, $Q=J+eZ$. Subtracting its square from (IC7) gives the
exact identity
\begin{align*}
 G-4DQ^2={}&(4BC-\kappa z/B)R_0^2+4BmR_1^2/z^2
       +2(2C+1)R_2^2\\
 &+r^2(4nw^2+4n+2)T^2
 +\frac HB\{4A(1+u^2)+4C+2\}E^2\\
 &+\kappa uR_0R_1/B+(-4r+\kappa/B)R_2T
       -4a(2C+1)R_2E.
\tag{WIE4}
\end{align*}
Add the exact zero (IC8a). Restore the full normals
\[
 E_X=\sqrt N(X+ruZ),\qquad
 E_Y=\sqrt{B/[H(1-H)]}(Y-HwZ).
\]
Use the exact relation
\[
 R_2=F_D/\sqrt D+rT+aE-uE_X/\sqrt N
                         +wE_Y/\sqrt{B/[H(1-H)]}.
\tag{WIE5}
\]
The inverse $1/\sqrt D$ is purely in the source field and infinitesimal.
All coefficients after this substitution are bounded.

For completeness, the first-normal quotients of
$\widehat F-4F_D^2$ are precisely the D-removed versions of
(IC11)--(IC12). Explicitly,
\begin{align*}
 \bar f_X&=-2nr(1+V)-r,\\
 \bar g_X&=-2mB/z-2nr(1+z)+1-H-r-2m(H+r-1),\\
 \bar f&=2m-2C(1+V)-V,\\
 \bar f_Y&=2AH(1+V)-2\alpha(1-H)+(2m+1)(H+r-1),\\
 \bar g_Y&=-B(2C+1)+2AH(1+z)+2CH+H.
\end{align*}
The constant term in the $Y$ monic polynomial acquires $2De$, but
its monic quotient coefficients are unchanged; the $X$ polynomial
has the constant correction $-2De$. Exact source stationarity
still sets the constant remainders to zero because $J+eZ$ vanishes
at the contact. Thus the displayed quotient formulas reconstruct
both linear normals. Their normalized coefficients are infinitesimal
by (WIE3); in particular $De$ is bounded and no bare $D$ survives.
The normal quadratic block of the D-removed expression has residue
$\operatorname{diag}(1,1)$.

In the order $(R_0,R_1,T,E,F_D,E_X,E_Y)$, the complete residue Gram
is diagonal:
\[
 \operatorname{diag}\{1,1-r_0,1+r_0-r_0^2,1,4,1,1\}.
\tag{WIE6}
\]
For instance the $T^2$ entry comes from substituting $R_2=r_0T$
in (WIE4): it is $r_0\nu_0-r_0^2$, plus the unit contribution from
(IC8a). Every entry is positive for $0\le r_0\le1/2$.
This proves full-parent tensor SOS at the corner also when
$1-\rho$ is infinitesimal. Along with the preserved interior proof
and the $\rho_0=0$ argument, every allowed source shape is covered
at the independent testing corner $u,w$ infinitesimal.
